% =====================================================================
%  Appendix E: Differential Forms Quick Reference
% =====================================================================

\chapter{Differential Forms Quick Reference}
\label{app:differential-forms-reference}

Differential forms are integrands with orientation built in.  A \(1\)-form is built for
curves.  A \(2\)-form is built for surfaces.  A \(3\)-form is built for solids.

This appendix is a reference.  It is meant to be scanned.

\[
\begin{array}{c|c|c}
\text{Form} & \text{Integrated over} & \text{Typical meaning} \\ \hline
0\text{-form} & \text{points} & \text{function value}\\[1mm]
1\text{-form} & \text{curves} & \text{work, circulation}\\[1mm]
2\text{-form} & \text{surfaces} & \text{flux, signed area}\\[1mm]
3\text{-form} & \text{solids} & \text{density, signed volume}
\end{array}
\]

% ---------------------------------------------------------------------
\section{\texorpdfstring{\(0\)}{0}-forms, \texorpdfstring{\(1\)}{1}-forms,
\texorpdfstring{\(2\)}{2}-forms, and \texorpdfstring{\(3\)}{3}-forms}
\label{app:E.1}

\subsection*{\texorpdfstring{\(0\)}{0}-forms}

A \(0\)-form is a function:
\[
\boxed{
    f.
}
\]

Examples:
\[
    f(x),
    \qquad
    f(x,y),
    \qquad
    f(x,y,z).
\]

A \(0\)-form is integrated over signed points:
\[
\boxed{
    \int_{+P}f=f(P),
    \qquad
    \int_{-P}f=-f(P).
}
\]

For an oriented interval,
\[
\boxed{
    \partial[a,b]=\{b\}-\{a\}.
}
\]

So
\[
\boxed{
    \int_{\partial[a,b]}f=f(b)-f(a).
}
\]

\subsection*{\texorpdfstring{\(1\)}{1}-forms}

A \(1\)-form in the plane has the form
\[
\boxed{
    \omega=P(x,y)\,dx+Q(x,y)\,dy.
}
\]

A \(1\)-form in space has the form
\[
\boxed{
    \omega=P(x,y,z)\,dx+Q(x,y,z)\,dy+R(x,y,z)\,dz.
}
\]

It is integrated over oriented curves.

If
\[
    \mathbf F=(P,Q,R),
\]
then the work \(1\)-form is
\[
\boxed{
    \omega_{\mathbf F}=P\,dx+Q\,dy+R\,dz.
}
\]

Line integral:
\[
\boxed{
    \int_C \mathbf F\cdot d\mathbf r
    =
    \int_C \omega_{\mathbf F}.
}
\]

\subsection*{\texorpdfstring{\(2\)}{2}-forms in the plane}

A \(2\)-form in the plane has the form
\[
\boxed{
    \alpha=f(x,y)\,dx\wedge dy.
}
\]

It is integrated over oriented plane regions.

The form
\[
\boxed{
    dx\wedge dy
}
\]
measures signed area.

\[
\boxed{
    \iint_R f\,dx\wedge dy
}
\]
means signed accumulation of \(f\) over \(R\).

For ordinary unsigned area,
\[
\boxed{
    dA=|dx\wedge dy|
}
\]
in the usual orientation.  Usually we write
\[
\boxed{
    dA=dx\,dy
}
\]
for positive area.

\subsection*{\texorpdfstring{\(2\)}{2}-forms in space}

A general \(2\)-form in space can be written as
\[
\boxed{
    \Phi
    =
    A\,dy\wedge dz
    +
    B\,dz\wedge dx
    +
    C\,dx\wedge dy.
}
\]

This ordering matches flux.

If
\[
    \mathbf F=(P,Q,R),
\]
the flux \(2\)-form is
\[
\boxed{
    \Phi_{\mathbf F}
    =
    P\,dy\wedge dz
    +
    Q\,dz\wedge dx
    +
    R\,dx\wedge dy.
}
\]

Flux integral:
\[
\boxed{
    \iint_S \mathbf F\cdot\widehat{\mathbf n}\,dS
    =
    \iint_S \Phi_{\mathbf F}.
}
\]

\subsection*{\texorpdfstring{\(3\)}{3}-forms}

A \(3\)-form in space has the form
\[
\boxed{
    \Omega=f(x,y,z)\,dx\wedge dy\wedge dz.
}
\]

It is integrated over oriented solids.

The form
\[
\boxed{
    dx\wedge dy\wedge dz
}
\]
measures signed volume.

For ordinary volume,
\[
\boxed{
    dV=dx\,dy\,dz
}
\]
in the usual orientation.

\subsection*{Degree}

The degree of a form is the number of differentials wedged together.

\[
\begin{array}{c|c}
\text{Degree} & \text{Example} \\ \hline
0 & f\\
1 & P\,dx+Q\,dy\\
2 & f\,dx\wedge dy\\
3 & f\,dx\wedge dy\wedge dz
\end{array}
\]

\subsection*{Common warnings}

A \(1\)-form is not a vector field, though a vector field can produce a \(1\)-form.

A \(2\)-form is not a vector field, though a vector field can produce a flux \(2\)-form.

Forms remember orientation.  Reversing orientation changes the sign of the integral.

% ---------------------------------------------------------------------
\section{Wedge product rules}
\label{app:E.2}

\subsection*{Basic antisymmetry}

\[
\boxed{
    dx\wedge dy=-dy\wedge dx
}
\]

\[
\boxed{
    dy\wedge dz=-dz\wedge dy
}
\]

\[
\boxed{
    dz\wedge dx=-dx\wedge dz
}
\]

Swapping two factors changes the sign.

\subsection*{Repeated factors vanish}

\[
\boxed{
    dx\wedge dx=0
}
\]

\[
\boxed{
    dy\wedge dy=0
}
\]

\[
\boxed{
    dz\wedge dz=0
}
\]

Any wedge product with a repeated differential is zero.

\subsection*{Positive volume order}

\[
\boxed{
    dx\wedge dy\wedge dz
}
\]
is the positive volume form in the usual orientation of \(\mathbb R^3\).

Cyclic orders are positive:
\[
\boxed{
    dx\wedge dy\wedge dz
    =
    dy\wedge dz\wedge dx
    =
    dz\wedge dx\wedge dy.
}
\]

Reversed cyclic orders are negative:
\[
\boxed{
    dx\wedge dz\wedge dy
    =
    dy\wedge dx\wedge dz
    =
    dz\wedge dy\wedge dx
    =
    -dx\wedge dy\wedge dz.
}
\]

\subsection*{Distributive rules}

\[
\boxed{
    \alpha\wedge(\beta+\gamma)
    =
    \alpha\wedge\beta+\alpha\wedge\gamma
}
\]

\[
\boxed{
    (\alpha+\beta)\wedge\gamma
    =
    \alpha\wedge\gamma+\beta\wedge\gamma
}
\]

Functions can be moved through wedge products:
\[
\boxed{
    (f\alpha)\wedge\beta
    =
    f(\alpha\wedge\beta).
}
\]

\subsection*{Graded sign rule}

If \(\alpha\) is a \(k\)-form and \(\beta\) is an \(\ell\)-form, then
\[
\boxed{
    \alpha\wedge\beta
    =
    (-1)^{k\ell}\beta\wedge\alpha.
}
\]

Examples:
\[
\boxed{
    \text{\(1\)-forms anticommute.}
}
\]

\[
\boxed{
    \text{A \(1\)-form and a \(2\)-form commute.}
}
\]

\subsection*{Two \texorpdfstring{\(1\)}{1}-forms in the plane}

Let
\[
    \alpha=a\,dx+b\,dy,
    \qquad
    \beta=c\,dx+d\,dy.
\]

Then
\[
\boxed{
    \alpha\wedge\beta
    =
    (ad-bc)\,dx\wedge dy.
}
\]

This is the determinant:
\[
\boxed{
    ad-bc=
    \det
    \begin{pmatrix}
    a&c\\
    b&d
    \end{pmatrix}.
}
\]

\subsection*{Two \texorpdfstring{\(1\)}{1}-forms in space}

Let
\[
    \alpha=a_1\,dx+a_2\,dy+a_3\,dz,
\]
and
\[
    \beta=b_1\,dx+b_2\,dy+b_3\,dz.
\]

Then
\[
\boxed{
\begin{aligned}
\alpha\wedge\beta
&=
(a_2b_3-a_3b_2)\,dy\wedge dz\\
&\quad+
(a_3b_1-a_1b_3)\,dz\wedge dx\\
&\quad+
(a_1b_2-a_2b_1)\,dx\wedge dy.
\end{aligned}
}
\]

The coefficients are the components of
\[
\boxed{
    \mathbf a\times\mathbf b.
}
\]

\subsection*{Three \texorpdfstring{\(1\)}{1}-forms in space}

Let
\[
    \alpha=a_1dx+a_2dy+a_3dz,
\]
\[
    \beta=b_1dx+b_2dy+b_3dz,
\]
\[
    \gamma=c_1dx+c_2dy+c_3dz.
\]

Then
\[
\boxed{
    \alpha\wedge\beta\wedge\gamma
    =
    \det
    \begin{pmatrix}
    a_1&b_1&c_1\\
    a_2&b_2&c_2\\
    a_3&b_3&c_3
    \end{pmatrix}
    dx\wedge dy\wedge dz.
}
\]

\subsection*{Common warnings}

The wedge product is not ordinary multiplication.

Order matters.

Repeated differentials vanish.

In \(\mathbb R^3\), every \(4\)-form is zero because there are only three coordinate
directions.

% ---------------------------------------------------------------------
\section{Pullbacks}
\label{app:E.3}

\subsection*{Meaning}

A pullback rewrites a form in new parameters.

If
\[
    F:\text{parameter space}\to\text{physical space},
\]
then
\[
\boxed{
    F^*\omega
}
\]
is the form \(\omega\) written in the parameter variables.

Pullbacks are the form version of substitution.

\subsection*{Pullback along a curve}

Let
\[
    \mathbf r(t)=(x(t),y(t),z(t)).
\]

Then
\[
\boxed{
    \mathbf r^*(dx)=x'(t)\,dt,
}
\]

\[
\boxed{
    \mathbf r^*(dy)=y'(t)\,dt,
}
\]

\[
\boxed{
    \mathbf r^*(dz)=z'(t)\,dt.
}
\]

For
\[
    \omega=P\,dx+Q\,dy+R\,dz,
\]
\[
\boxed{
    \mathbf r^*\omega
    =
    \bigl(P(\mathbf r(t))x'(t)
    +
    Q(\mathbf r(t))y'(t)
    +
    R(\mathbf r(t))z'(t)\bigr)\,dt.
}
\]

Therefore
\[
\boxed{
    \int_C \omega
    =
    \int_a^b \mathbf r^*\omega.
}
\]

\subsection*{Pullback over a parametrized surface}

Let
\[
    \mathbf r(u,v)=(x(u,v),y(u,v),z(u,v)).
\]

Then
\[
\boxed{
    \mathbf r^*(dx)=x_u\,du+x_v\,dv.
}
\]

\[
\boxed{
    \mathbf r^*(dy)=y_u\,du+y_v\,dv.
}
\]

\[
\boxed{
    \mathbf r^*(dz)=z_u\,du+z_v\,dv.
}
\]

\subsection*{Pullback of coordinate area forms}

\[
\boxed{
    \mathbf r^*(dx\wedge dy)
    =
    (x_u y_v-x_v y_u)\,du\wedge dv.
}
\]

\[
\boxed{
    \mathbf r^*(dy\wedge dz)
    =
    (y_u z_v-y_v z_u)\,du\wedge dv.
}
\]

\[
\boxed{
    \mathbf r^*(dz\wedge dx)
    =
    (z_u x_v-z_v x_u)\,du\wedge dv.
}
\]

\subsection*{Flux form pullback}

Let
\[
    \Phi_{\mathbf F}
    =
    P\,dy\wedge dz
    +
    Q\,dz\wedge dx
    +
    R\,dx\wedge dy.
\]

Then
\[
\boxed{
    \mathbf r^*\Phi_{\mathbf F}
    =
    \mathbf F(\mathbf r(u,v))
    \cdot
    (\mathbf r_u\times\mathbf r_v)
    \,du\wedge dv.
}
\]

Thus
\[
\boxed{
    \iint_S \Phi_{\mathbf F}
    =
    \iint_D
    \mathbf F(\mathbf r(u,v))
    \cdot
    (\mathbf r_u\times\mathbf r_v)
    \,du\,dv.
}
\]

\subsection*{Pullback under a coordinate change}

Let
\[
    T(u,v)=(x(u,v),y(u,v)).
\]

Then
\[
\boxed{
    T^*(dx\wedge dy)
    =
    \frac{\partial(x,y)}{\partial(u,v)}
    \,du\wedge dv.
}
\]

That is,
\[
\boxed{
    dx\wedge dy
    =
    \det
    \begin{pmatrix}
    x_u&x_v\\
    y_u&y_v
    \end{pmatrix}
    du\wedge dv.
}
\]

For a space coordinate change
\[
    T(u,v,w)=(x,y,z),
\]
\[
\boxed{
    T^*(dx\wedge dy\wedge dz)
    =
    \frac{\partial(x,y,z)}{\partial(u,v,w)}
    \,du\wedge dv\wedge dw.
}
\]

\subsection*{Polar pullback}

For
\[
    x=r\cos\theta,
    \qquad
    y=r\sin\theta,
\]
\[
\boxed{
    dx\wedge dy
    =
    r\,dr\wedge d\theta.
}
\]

\subsection*{Cylindrical pullback}

For
\[
    x=r\cos\theta,
    \qquad
    y=r\sin\theta,
    \qquad
    z=z,
\]
\[
\boxed{
    dx\wedge dy\wedge dz
    =
    r\,dr\wedge d\theta\wedge dz.
}
\]

\subsection*{Spherical pullback}

For
\[
    x=\rho\sin\phi\cos\theta,
\]
\[
    y=\rho\sin\phi\sin\theta,
\]
\[
    z=\rho\cos\phi,
\]
\[
\boxed{
    dx\wedge dy\wedge dz
    =
    \rho^2\sin\phi\,
    d\rho\wedge d\phi\wedge d\theta.
}
\]

\subsection*{Common warnings}

Pullbacks respect wedge products:
\[
\boxed{
    F^*(\alpha\wedge\beta)=F^*\alpha\wedge F^*\beta.
}
\]

Pullbacks respect exterior derivatives:
\[
\boxed{
    F^*(d\omega)=d(F^*\omega).
}
\]

Ordinary change of variables uses absolute values for positive area and volume.

Differential forms keep orientation signs.

% ---------------------------------------------------------------------
\section{Exterior derivative formulas}
\label{app:E.4}

\subsection*{Main idea}

The exterior derivative raises degree by one:
\[
\boxed{
    d:\text{\(k\)-forms}\longrightarrow \text{\((k+1)\)-forms}.
}
\]

\[
\begin{array}{c|c}
\text{Input} & \text{Output} \\ \hline
0\text{-form} & 1\text{-form}\\
1\text{-form} & 2\text{-form}\\
2\text{-form} & 3\text{-form}\\
3\text{-form} & 0\text{ in }\mathbb R^3
\end{array}
\]

\subsection*{Derivative of a \texorpdfstring{\(0\)}{0}-form}

In the plane:
\[
\boxed{
    df=f_x\,dx+f_y\,dy.
}
\]

In space:
\[
\boxed{
    df=f_x\,dx+f_y\,dy+f_z\,dz.
}
\]

This corresponds to the gradient.

\subsection*{Derivative of a plane \texorpdfstring{\(1\)}{1}-form}

Let
\[
    \omega=P\,dx+Q\,dy.
\]

Then
\[
\boxed{
    d\omega
    =
    (Q_x-P_y)\,dx\wedge dy.
}
\]

This is scalar curl times signed area.

\subsection*{Derivative of a space \texorpdfstring{\(1\)}{1}-form}

Let
\[
    \omega=P\,dx+Q\,dy+R\,dz.
\]

Then
\[
\boxed{
\begin{aligned}
d\omega
&=
(R_y-Q_z)\,dy\wedge dz\\
&\quad+
(P_z-R_x)\,dz\wedge dx\\
&\quad+
(Q_x-P_y)\,dx\wedge dy.
\end{aligned}
}
\]

This corresponds to curl.

\subsection*{Derivative of a space \texorpdfstring{\(2\)}{2}-form}

Let
\[
    \Phi=P\,dy\wedge dz+Q\,dz\wedge dx+R\,dx\wedge dy.
\]

Then
\[
\boxed{
    d\Phi
    =
    (P_x+Q_y+R_z)\,dx\wedge dy\wedge dz.
}
\]

This corresponds to divergence.

\subsection*{Derivative of a \texorpdfstring{\(3\)}{3}-form in \(\mathbb R^3\)}

If
\[
    \Omega=f\,dx\wedge dy\wedge dz,
\]
then
\[
\boxed{
    d\Omega=0
}
\]
in \(\mathbb R^3\).

There are no nonzero \(4\)-forms in three-dimensional space.

\subsection*{Product rule}

If \(\alpha\) is a \(k\)-form, then
\[
\boxed{
    d(\alpha\wedge\beta)
    =
    d\alpha\wedge\beta
    +
    (-1)^k\alpha\wedge d\beta.
}
\]

Special case for a function \(f\):
\[
\boxed{
    d(f\alpha)=df\wedge\alpha+f\,d\alpha.
}
\]

\subsection*{Exterior derivative squared}

\[
\boxed{
    d(d\omega)=0.
}
\]

Usually written:
\[
\boxed{
    d^2=0.
}
\]

Classical meanings:
\[
\boxed{
    \operatorname{curl}(\nabla f)=\mathbf 0.
}
\]

\[
\boxed{
    \operatorname{div}(\operatorname{curl}\mathbf F)=0.
}
\]

\subsection*{Closed and exact forms}

A form \(\omega\) is closed if
\[
\boxed{
    d\omega=0.
}
\]

A form \(\omega\) is exact if
\[
\boxed{
    \omega=d\eta
}
\]
for some form \(\eta\).

Exact forms are always closed:
\[
\boxed{
    \omega=d\eta
    \quad\Longrightarrow\quad
    d\omega=d^2\eta=0.
}
\]

The converse is not always true.  Holes in the domain can prevent a closed form from
being exact.

\subsection*{Common warnings}

Partial derivative notation depends on the coordinate variables being used.

The exterior derivative does not depend on a metric.

The gradient, curl, and divergence interpretations use Euclidean structure to translate
forms back into vectors.

% ---------------------------------------------------------------------
\section{Classical vector calculus translated into forms}
\label{app:E.5}

\subsection*{Gradient}

Classical:
\[
\boxed{
    \nabla f=(f_x,f_y,f_z).
}
\]

Forms:
\[
\boxed{
    df=f_x\,dx+f_y\,dy+f_z\,dz.
}
\]

Translation:
\[
\boxed{
    df(\mathbf v)=\nabla f\cdot \mathbf v.
}
\]

\subsection*{Work form}

For
\[
    \mathbf F=(P,Q,R),
\]
classical line integral:
\[
\boxed{
    \int_C \mathbf F\cdot d\mathbf r.
}
\]

Forms:
\[
\boxed{
    \omega_{\mathbf F}=P\,dx+Q\,dy+R\,dz.
}
\]

\[
\boxed{
    \int_C \mathbf F\cdot d\mathbf r
    =
    \int_C \omega_{\mathbf F}.
}
\]

\subsection*{Curl}

Classical:
\[
\boxed{
    \nabla\times\mathbf F
    =
    (R_y-Q_z,\ P_z-R_x,\ Q_x-P_y).
}
\]

Forms:
\[
\boxed{
    d\omega_{\mathbf F}
    =
    (R_y-Q_z)\,dy\wedge dz
    +
    (P_z-R_x)\,dz\wedge dx
    +
    (Q_x-P_y)\,dx\wedge dy.
}
\]

Translation:
\[
\boxed{
    d\omega_{\mathbf F}
    =
    \Phi_{\nabla\times\mathbf F}.
}
\]

\subsection*{Flux form}

For
\[
    \mathbf F=(P,Q,R),
\]
classical flux:
\[
\boxed{
    \iint_S \mathbf F\cdot\widehat{\mathbf n}\,dS.
}
\]

Forms:
\[
\boxed{
    \Phi_{\mathbf F}
    =
    P\,dy\wedge dz
    +
    Q\,dz\wedge dx
    +
    R\,dx\wedge dy.
}
\]

\[
\boxed{
    \iint_S \mathbf F\cdot\widehat{\mathbf n}\,dS
    =
    \iint_S \Phi_{\mathbf F}.
}
\]

\subsection*{Divergence}

Classical:
\[
\boxed{
    \nabla\cdot\mathbf F=P_x+Q_y+R_z.
}
\]

Forms:
\[
\boxed{
    d\Phi_{\mathbf F}
    =
    (P_x+Q_y+R_z)\,dx\wedge dy\wedge dz.
}
\]

Translation:
\[
\boxed{
    d\Phi_{\mathbf F}
    =
    (\nabla\cdot\mathbf F)\,dV.
}
\]

\subsection*{Plane circulation}

For
\[
    \mathbf F=(P,Q),
\]
work form:
\[
\boxed{
    \omega=P\,dx+Q\,dy.
}
\]

Scalar curl:
\[
\boxed{
    \operatorname{curl}_{\mathrm{plane}}\mathbf F=Q_x-P_y.
}
\]

Forms:
\[
\boxed{
    d\omega=(Q_x-P_y)\,dx\wedge dy.
}
\]

\subsection*{Plane flux}

For
\[
    \mathbf F=(M,N),
\]
the outward flux form on a positively oriented boundary is
\[
\boxed{
    \phi_{\mathbf F}=M\,dy-N\,dx.
}
\]

Then
\[
\boxed{
    \int_{\partial R}\phi_{\mathbf F}
    =
    \int_{\partial R}\mathbf F\cdot\widehat{\mathbf n}\,ds.
}
\]

Its exterior derivative is
\[
\boxed{
    d\phi_{\mathbf F}
    =
    (M_x+N_y)\,dx\wedge dy.
}
\]

So
\[
\boxed{
    d\phi_{\mathbf F}
    =
    (\nabla\cdot\mathbf F)\,dA.
}
\]

\subsection*{Vector identities from \texorpdfstring{\(d^2=0\)}{d squared equals zero}}

From
\[
    d(df)=0,
\]
we get
\[
\boxed{
    \nabla\times(\nabla f)=\mathbf 0.
}
\]

From
\[
    d(d\omega_{\mathbf F})=0,
\]
we get
\[
\boxed{
    \nabla\cdot(\nabla\times\mathbf F)=0.
}
\]

\subsection*{Common dictionary}

\[
\begin{array}{c|c}
\text{Classical language} & \text{Forms language} \\ \hline
f & 0\text{-form}\\[1mm]
\nabla f & df\\[1mm]
\mathbf F\cdot d\mathbf r & P\,dx+Q\,dy+R\,dz\\[1mm]
\nabla\times\mathbf F & d(P\,dx+Q\,dy+R\,dz)\\[1mm]
\mathbf F\cdot\widehat{\mathbf n}\,dS &
P\,dy\wedge dz+Q\,dz\wedge dx+R\,dx\wedge dy\\[1mm]
\nabla\cdot\mathbf F & d\Phi_{\mathbf F}
\end{array}
\]

\subsection*{Common warnings}

The same vector field \(\mathbf F\) can produce two different forms:

\[
\boxed{
    \omega_{\mathbf F}=P\,dx+Q\,dy+R\,dz
}
\]
for work and circulation.

\[
\boxed{
    \Phi_{\mathbf F}
    =
    P\,dy\wedge dz+Q\,dz\wedge dx+R\,dx\wedge dy
}
\]
for flux.

Do not confuse these.

% ---------------------------------------------------------------------
\section{Generalized Stokes theorem summary}
\label{app:E.6}

\subsection*{The theorem}

\[
\boxed{
    \int_{\partial M}\omega
    =
    \int_M d\omega.
}
\]

This is the generalized Stokes theorem.

It says:

\[
\boxed{
    \text{the integral of a form over the boundary}
    =
    \text{the integral of its derivative over the region}.
}
\]

\subsection*{Degree matching}

If \(M\) is \(k\)-dimensional, then \(\omega\) is a \((k-1)\)-form and \(d\omega\) is a
\(k\)-form.

\[
\begin{array}{c|c|c}
M & \omega & d\omega \\ \hline
\text{interval} & 0\text{-form} & 1\text{-form}\\[1mm]
\text{plane region} & 1\text{-form} & 2\text{-form}\\[1mm]
\text{surface} & 1\text{-form} & 2\text{-form}\\[1mm]
\text{solid} & 2\text{-form} & 3\text{-form}
\end{array}
\]

\subsection*{Fundamental Theorem of Calculus}

Let \(M=[a,b]\), and let \(\omega=f\).  Then
\[
    d\omega=df=f'(x)\,dx.
\]

Generalized Stokes gives
\[
\boxed{
    \int_a^b f'(x)\,dx=f(b)-f(a).
}
\]

Boundary:
\[
\boxed{
    \partial[a,b]=\{b\}-\{a\}.
}
\]

\subsection*{Green's theorem: circulation form}

Let
\[
    \omega=P\,dx+Q\,dy.
\]

Then
\[
    d\omega=(Q_x-P_y)\,dx\wedge dy.
\]

Stokes gives
\[
\boxed{
    \int_{\partial R}P\,dx+Q\,dy
    =
    \iint_R (Q_x-P_y)\,dx\wedge dy.
}
\]

Classical notation:
\[
\boxed{
    \oint_{\partial R}\mathbf F\cdot d\mathbf r
    =
    \iint_R \operatorname{curl}_{\mathrm{plane}}\mathbf F\,dA.
}
\]

\subsection*{Green's theorem: flux form}

Let
\[
    \mathbf F=(M,N),
\]
and let
\[
    \phi_{\mathbf F}=M\,dy-N\,dx.
\]

Then
\[
    d\phi_{\mathbf F}=(M_x+N_y)\,dx\wedge dy.
\]

Stokes gives
\[
\boxed{
    \int_{\partial R}M\,dy-N\,dx
    =
    \iint_R (M_x+N_y)\,dx\wedge dy.
}
\]

Classical notation:
\[
\boxed{
    \int_{\partial R}\mathbf F\cdot\widehat{\mathbf n}\,ds
    =
    \iint_R \nabla\cdot\mathbf F\,dA.
}
\]

\subsection*{Stokes' theorem in space}

Let
\[
    \omega=P\,dx+Q\,dy+R\,dz.
\]

Then
\[
    d\omega=\Phi_{\nabla\times\mathbf F}.
\]

Stokes gives
\[
\boxed{
    \int_{\partial S}\omega
    =
    \int_S d\omega.
}
\]

Classical notation:
\[
\boxed{
    \oint_{\partial S}\mathbf F\cdot d\mathbf r
    =
    \iint_S(\nabla\times\mathbf F)\cdot\widehat{\mathbf n}\,dS.
}
\]

\subsection*{Divergence theorem}

Let
\[
    \Phi_{\mathbf F}
    =
    P\,dy\wedge dz+Q\,dz\wedge dx+R\,dx\wedge dy.
\]

Then
\[
    d\Phi_{\mathbf F}
    =
    (P_x+Q_y+R_z)\,dx\wedge dy\wedge dz.
\]

Stokes gives
\[
\boxed{
    \int_{\partial E}\Phi_{\mathbf F}
    =
    \int_E d\Phi_{\mathbf F}.
}
\]

Classical notation:
\[
\boxed{
    \iint_{\partial E}\mathbf F\cdot\widehat{\mathbf n}\,dS
    =
    \iiint_E\nabla\cdot\mathbf F\,dV.
}
\]

\subsection*{Orientation rules}

Interval:
\[
\boxed{
    \partial[a,b]=\{b\}-\{a\}.
}
\]

Plane region:
\[
\boxed{
    \text{walk so the region stays on your left.}
}
\]

Region with a hole:
\[
\boxed{
    \text{outer boundary counterclockwise, inner boundary clockwise.}
}
\]

Surface:
\[
\boxed{
    \text{use the right-hand rule from the chosen normal.}
}
\]

Solid:
\[
\boxed{
    \text{boundary is oriented outward.}
}
\]

\subsection*{Hypotheses checklist}

Before using Stokes' theorem, check:

\[
\begin{array}{ll}
1. & \text{The form is defined on the region and its boundary.}\\[1mm]
2. & \text{The needed partial derivatives exist and are continuous.}\\[1mm]
3. & \text{The region is oriented.}\\[1mm]
4. & \text{The boundary has the induced orientation.}\\[1mm]
5. & \text{Singularities are not inside the region unless handled separately.}\\[1mm]
6. & \text{The surface or region is piecewise smooth.}
\end{array}
\]

\subsection*{Closed versus exact warning}

If
\[
    \omega=d\eta,
\]
then
\[
    d\omega=0.
\]

So exact implies closed.

But closed does not always imply exact.

A missing point or hole in the domain can prevent a potential from existing globally.

\subsection*{Boundary of a boundary}

A boundary has no boundary:
\[
\boxed{
    \partial(\partial M)=0.
}
\]

The form version is
\[
\boxed{
    d^2=0.
}
\]

These are two sides of the same idea.

% ---------------------------------------------------------------------
\section*{One-page formula summary}

\[
\boxed{
    \omega=P\,dx+Q\,dy+R\,dz
}
\]

\[
\boxed{
    \Phi_{\mathbf F}
    =
    P\,dy\wedge dz+Q\,dz\wedge dx+R\,dx\wedge dy
}
\]

\[
\boxed{
    df=f_x\,dx+f_y\,dy+f_z\,dz
}
\]

\[
\boxed{
    d(P\,dx+Q\,dy)
    =
    (Q_x-P_y)\,dx\wedge dy
}
\]

\[
\boxed{
\begin{aligned}
d(P\,dx+Q\,dy+R\,dz)
&=
(R_y-Q_z)\,dy\wedge dz\\
&\quad+
(P_z-R_x)\,dz\wedge dx\\
&\quad+
(Q_x-P_y)\,dx\wedge dy
\end{aligned}
}
\]

\[
\boxed{
    d(P\,dy\wedge dz+Q\,dz\wedge dx+R\,dx\wedge dy)
    =
    (P_x+Q_y+R_z)\,dx\wedge dy\wedge dz
}
\]

\[
\boxed{
    dx\wedge dy=-dy\wedge dx
}
\]

\[
\boxed{
    dx\wedge dx=0
}
\]

\[
\boxed{
    dx\wedge dy\wedge dz
    =
    dy\wedge dz\wedge dx
    =
    dz\wedge dx\wedge dy
}
\]

\[
\boxed{
    \mathbf r^*(P\,dx+Q\,dy+R\,dz)
    =
    (P x'+Q y'+R z')\,dt
}
\]

\[
\boxed{
    \mathbf r^*\Phi_{\mathbf F}
    =
    \mathbf F(\mathbf r(u,v))\cdot(\mathbf r_u\times\mathbf r_v)
    \,du\wedge dv
}
\]

\[
\boxed{
    dx\wedge dy=r\,dr\wedge d\theta
}
\]

\[
\boxed{
    dx\wedge dy\wedge dz
    =
    \rho^2\sin\phi\,
    d\rho\wedge d\phi\wedge d\theta
}
\]

\[
\boxed{
    d^2=0
}
\]

\[
\boxed{
    \int_{\partial M}\omega=\int_M d\omega
}
\]